\documentclass[11pt]{article}
\usepackage[a4paper,margin=24mm]{geometry}
\usepackage{amsmath,amssymb,amsthm,booktabs,mathtools}
\usepackage[hidelinks]{hyperref}
\allowdisplaybreaks[2]
\newtheorem{theorem}{Theorem}
\newtheorem{corollary}{Corollary}
\newcommand{\eps}{\varepsilon}
\begin{document}
\title{A fixed certificate for reciprocal Alice angles}
\author{}\date{23 October 2026}\maketitle

The source ratio can become large while the corresponding state quadratic
remains positive. Two source-dependent rescalings reveal why: the remaining
coupling vector has norm at most one. We use this contraction to complete
the leading quadratic form, then control the three corrections in the full
parent. The resulting testing functions are independent of the source.

\paragraph{The source and the polynomial to be certified.}
The source variables obey $x,\sigma,z>0$, $0<b,h<1$,
$r=bz<1$. Set
\begin{gather*}
 Q=1-b,\quad \theta=1-h,\quad e=xh/\sigma,\quad E_s=1-e,
 \quad K_s=1+x,\quad R_s=1+xr,\\
 M=R_s-xz\theta,\quad J=1+xh-\sigma(1-r),\\
 P_s=\sigma+K_s+\sigma z(1+bx+2xh)
                   +\sigma bx^2hz^2+x^2h\theta z,\\
 \pi=\frac{z(1+xh)E_s+(1-r)(1+xhz)}{2P_s},\qquad
 \ell=\frac{(1+\sigma r)\{eR_s(1+\sigma z)+ME_s\}}{2P_s},\\
 \kappa_s=\frac{K_s(1+\sigma r)(1+\sigma z)}{2MP_s}.
\end{gather*}
The strict gaps are $E_s,M,J>0$. The seven source coefficients,
in the normalization $d=1$, are
\begin{gather*}
 m=\kappa_s x^2z\frac h\sigma,\qquad
 C=\frac{K_s z(1+\sigma r)E_s}{2(1-r)P_s},\qquad
 A=\frac\pi{bz},\qquad \alpha=\frac\ell{Qz},\\
 D=\frac{K_s(1+\sigma z)(1+\sigma r)-JP_s}{2JP_s},\qquad
 \beta=\frac{K_s(1+\sigma r)(1+xhz)-\theta P_s}{2\theta P_s}.
\end{gather*}
In particular $m>0$; the hypotheses also require $D,\beta\ge0$.

For $D_p=1+p^2$, $D_w=1+w^2$, the complete inverse-Alice parent is
\begin{equation}\label{eq:allz-explicit-parent}
\begin{split}
 \widehat F={}&4m(pW-pwX+Y-wZ)^2+4CD_w(W-pY)^2
                         +4AD_pD_wY^2\\
 &+(4D+2)(wW+X-pwY-pZ)^2
        +(4\alpha+2)D_p(Y-wZ)^2\\
 &+(4\beta+2)D_w(pX+Z)^2-2D_pD_wZ^2.
\end{split}
\end{equation}
It equals $p^2F_w(1/p,w)$ for $p\ne0$. For unit states the physical
slack is $\widehat F/(D_pD_w)$. The pure testing factor
$\kappa_t=(D_pD_w)^{-1/2}$ therefore restores the physical normalization
and remains regular at $p=0$.



\begin{theorem}\label{thm:allz}
Fix a real constant $P_*\ge1$ and set
\[
 \eta=2^{-30}P_*^{-2},\qquad
 \delta=2^{-12}P_*^{-1},\qquad f=1-25\eta.
\]
Use the native source just defined, normalized by $d=1$. Suppose
\begin{equation}\label{eq:domain}
 0<x,\sigma\le\frac\eta{1+z},\quad z>0,\quad
 0<b,h<1,\quad r=bz<1,\quad E_s,M,J>0,\quad D,\beta\ge0.
\end{equation}
There is no upper bound on $H=h/\sigma$. For the independent testing
coordinates
\[
 p=1/u,\qquad 0\le p\le P_*,\quad 0\le w\le\delta,
\]
the complete inverse-Alice parent has a separated square identity in
one fixed finite family of testing functions. This family depends on
$P_*$ and $\delta$, and is independent of $z$ or any upper bound for
$z$. The construction uses $51$ testing functions. The identity
includes every state vector and the literal testing face $p=0$.
Neither $D$ nor $\beta$ is inverted, so their permitted zero faces
are included.
\end{theorem}

The testing statement applies to the projectively normalized parent
$p^2F(1/p,w)$, continued polynomially at $p=0$. Negative angle signs
are handled by the fixed gauges
$(u,Y,Z)\mapsto(-u,-Y,-Z)$ and
$(w,X,Z)\mapsto(-w,-X,-Z)$. Complex states are separated into real
and imaginary parts. No restriction on a ratio of state coordinates
is imposed.

\paragraph{Source estimates.}
Define
\begin{gather*}
 H_0=1+zQ,\qquad B_s=P_s+xz(J-\sigma),\\
 \mathcal H=R_sH_0+z\{(1-r)xh(2+xr)
                  -e[2+x(1+h+r\theta)]+K_s e^2\}.
\end{gather*}
The exact native identities are
\begin{equation}\label{eq:BE}
 2\theta P_s\beta=hB_s-\sigma H_0,\qquad
 2P_s(JD+x\theta\beta)=\sigma^2\mathcal H.
\end{equation}
Define the leading weights
\begin{equation}\label{eq:leadingweights}
 C_0=\frac{z}{2(1-r)},\qquad
 A_0=\frac{H_0}{2bz},\qquad
 \alpha_0=\frac1{2Qz}.
\end{equation}
The native coefficients give the exact ratios
\begin{align}
 \frac C{C_0}&=\frac{K_s(1+\sigma r)(1-e)}{P_s},\label{eq:ratioC}\\
 \frac A{A_0}&=
 \frac{z(1+xh)(1-e)+(1-r)(1+xhz)}{P_sH_0},\label{eq:ratioA}\\
 \frac\alpha{\alpha_0}&=
 \frac{(1+\sigma r)\{eR_s(1+\sigma z)+M(1-e)\}}{P_s}.
 \label{eq:ratioalpha}
\end{align}

To control all three ratios together, bound the weighted source openings.
Put $T=1+z$. From \eqref{eq:domain}, $Tx,T\sigma\le\eta$.
Using $bz<1$ in the $\sigma bx^2hz^2$ term shows
\[
 1\le P_s\le1+4\eta,\quad |M-1|\le\eta,\quad
 J-\sigma>0,\quad 1\le B_s\le1+6\eta.
\]
Initially $e<1$. Dropping the negative term in $\mathcal H$ yields
$\mathcal H\le4T$. From \eqref{eq:BE} one obtains the exact identity
\[
 eB_s=xH_0+\sigma\mathcal H-2P_sJD/\sigma,
\]
so $e\le Tx+4T\sigma\le5\eta$, uniformly in $H$. The first ratio lies between
$(1-5\eta)/(1+4\eta)\ge1-9\eta$ and
$(1+\eta)^2\le1+3\eta$.
The numerator of the second ratio lies between
$H_0(1-e)$ and $H_0+2xhz\le H_0(1+2\eta)$, since $H_0\ge z$.
It has the same lower bound $1-9\eta$ and upper bound $1+2\eta$.
For the third ratio the expression in braces is a convex combination
of $M$ and $R_s(1+\sigma z)$; they lie between $1-\eta$ and
$1+3\eta$. Consequently
\begin{equation}\label{eq:ratioest}
 \left|C/C_0-1\right|\le10\eta,\quad
 \left|A/A_0-1\right|\le10\eta,\quad
 \left|\alpha/\alpha_0-1\right|\le5\eta.
\end{equation}
Thus $C-fC_0$, $A-fA_0$, and $\alpha-f\alpha_0$ are nonnegative.
The same fraction $f$ of each leading weight can therefore be retained.

\paragraph{The exact leading completion.}
Set $q=\sqrt z>0$ and introduce the full-state forms
\[
 E=qW,\quad F=Y/q,\quad G=wZ/q,\quad
 R=X+wW,\quad S=pX.
\]
The complete core is
\begin{equation}\label{eq:core}
 K=R^2+(1+w^2)S^2+z(G-F)^2
    +p^2(1+w^2)zF^2
    -2pE(G+w^2F)+2qwS(G-F).
\end{equation}
Write
\[
 K_0=R^2+S^2+z(G-F)^2+p^2zF^2-2pEG.
\]
Taking the three weights \eqref{eq:leadingweights} at $w=0$ gives
\begin{align*}
 Q_0={}&2K_0+\frac{4C_0}{z}(E-pzF)^2
          +4A_0z(1+p^2)F^2
          +4\alpha_0z(1+p^2)(F-G)^2.
\end{align*}
Define
\begin{gather*}
 J_0=F-bG,\quad B=\frac{H_0}{bQ},\quad
 J_* =\sqrt B\,J_0,\quad G_* =\sqrt{H_0}\,G,\\
 a=\frac z{\sqrt B},\qquad c=\frac1{\sqrt{H_0}},\qquad
 L=aJ_*+cG_*=zJ_0+G.
\end{gather*}
Direct completion of $E$ gives the exact identity
\begin{equation}\label{eq:leadingidentity}
 Q_0=2R^2+2S^2+\frac2{1-r}(E-pL)^2
              +2(1+p^2)J_*^2+2G_*^2.
\end{equation}
Indeed the $p^2$ remainder after completing $E$ is
$(z+1/Q)(F-bG)^2/b$, and the constant remainder diagonalizes in
$(J_0,G)$. The bounded source coefficients satisfy
\begin{equation}\label{eq:ac}
 a^2+c^2=\frac{bQz^2+1}{H_0}\le1,
 \qquad H_0-1-bQz^2=zQ(1-r)\ge0.
\end{equation}
This is the reason that a large value of $z$ causes no loss after the
source-only normalization.

Put $t=p/P_*\in[0,1]$, $\bar E=E/P_*$, and
$y=(\bar E,J_*,G_*,R,S)$. Let
\[
 B_0=2R^2+2S^2+2(\bar E-tL)^2
                         +2(1+t^2)J_*^2+2G_*^2.
\]
The leading excess is exactly
\begin{equation}\label{eq:leadingexcess}
 Q_0-B_0=2\left(\frac{P_*^2}{1-r}-1\right)(\bar E-tL)^2
                   +2(P_*^2-1)t^2J_*^2.
\end{equation}
It is nonnegative. More specifically, the following square identity
proves $B_0\succeq I_5/2$ on the independent forms $y$:
\begin{align}\label{eq:strictbaseline}
 B_0-\tfrac12\|y\|^2={}&
 \tfrac32(R^2+S^2)+\tfrac32(\bar E-\tfrac43tL)^2
 +\tfrac32J_*^2+\tfrac56G_*^2\nonumber\\
 &+\tfrac23(1-t^2)G_*^2+\tfrac43t^2J_*^2\nonumber\\
 &+\tfrac23t^2\{(1-a^2-c^2)(J_*^2+G_*^2)
                            +(cJ_*-aG_*)^2\}.
\end{align}
Every coefficient or radical in this identity belongs entirely to
the source field or entirely to the testing field.

\paragraph{The full retained-square identity.}
Let $\widehat F=p^2F(1/p,w)$ be the projectively normalized parent.
Besides the core $2K$, its three selected events are
\[
 \frac{4C}{z}(1+w^2)(E-pzF)^2,\qquad
 4Az(1+p^2)(1+w^2)F^2,\qquad
 4\alpha z(1+p^2)(F-G)^2.
\]
All remaining events are retained literally:
\begin{align*}
 \mathcal X={}&4m(pW-pwX+Y-wZ)^2
 +4D(wW+X-pwY-pZ)^2\\
 &+4\beta(1+w^2)(pX+Z)^2\\
 &+\frac{4(C-fC_0)}z(1+w^2)(E-pzF)^2\\
 &+4(A-fA_0)z(1+p^2)(1+w^2)F^2\\
 &+4(\alpha-f\alpha_0)z(1+p^2)(F-G)^2.
\end{align*}
They have nonnegative source weights. Define the further retained sum
\begin{align*}
 \mathcal P={}&2(1-f)\{R^2+S^2+z(G-F)^2+p^2zF^2\}\\
 &+2w^2S^2+2p^2w^2zF^2\\
 &+\frac{4fC_0}z w^2(E-pzF)^2
                      +4fA_0z(1+p^2)w^2F^2.
\end{align*}
Then the complete parent, with no omitted state terms, is
\begin{equation}\label{eq:whole}
 \widehat F=fQ_0+\mathcal E+\mathcal P+\mathcal X,
\end{equation}
where the only correction is
\begin{equation}\label{eq:error}
 \mathcal E=-4(1-f)pEG-4pw^2EJ_0-4pbw^2EG
                              -4qwSJ_0+4qQwSG.
\end{equation}
This follows from \eqref{eq:core}; its difference from $K_0$ includes
the positive term $p^2w^2zF^2$, which has been retained in
$\mathcal P$.

\paragraph{Explicit source Gram assignments.}
In the ordering $y=(\bar E,J_*,G_*,R,S)$ write
\[
 \mathcal E=y^{\mathsf T}(tC_1+tw^2C_2+wC_3)y.
\]
Only the following upper-triangular entries are nonzero; the lower
triangles are obtained by symmetry:
\begin{align*}
 &(C_1)_{13}=-2(1-f)P_*^2/\sqrt{H_0},\\
 &(C_2)_{12}=-2P_*^2/\sqrt B,\qquad
   (C_2)_{13}=-2bP_*^2/\sqrt{H_0},\\
 &(C_3)_{25}=-2q/\sqrt B,\qquad
   (C_3)_{35}=2qQ/\sqrt{H_0}.
\end{align*}
The row norms are at most $50\eta P_*^2,4P_*^2,4$, respectively.
For the last bound use
\[
 q^2/B=rQ/H_0\le1,\qquad q^2Q^2/H_0=zQ^2/H_0\le1.
\]
The first two use $H_0\ge1$ and $B\ge1$.
Set
\begin{gather*}
 T_1=C_1,\quad T_2=\delta^2C_2,\quad T_3=\delta C_3,\\
 M_1=50\eta P_*^2,\quad M_2=4P_*^2\delta^2,\quad
 M_3=4\delta,\quad \rho=M_1+M_2+M_3.
\end{gather*}
Then
\[
 \rho\le50\cdot2^{-30}+4\cdot2^{-24}+4\cdot2^{-12}<1/512,
 \qquad f/2-\rho>1/4.
\]
All six source matrices $(M_iI_5\pm T_i)/2$ are positive
semidefinite. With pure testing monomials
\[
 m_1=t,\quad m_2=t(w/\delta)^2,\quad m_3=w/\delta,
\]
one has the exact identity
\begin{equation}\label{eq:pairs}
 \rho\|y\|^2+\mathcal E
 =\sum_{i=1}^3\left\{
 (1+m_i)y^{\mathsf T}\frac{M_iI_5+T_i}2y
 +(1-m_i)y^{\mathsf T}\frac{M_iI_5-T_i}2y\right\}.
\end{equation}
Also
\[
 fQ_0-\rho\|y\|^2
 =f(Q_0-B_0)+f(B_0-\tfrac12\|y\|^2)
                     +(f/2-\rho)\|y\|^2.
\]
Equations \eqref{eq:leadingexcess}, \eqref{eq:strictbaseline}, and
\eqref{eq:pairs}, together with $\mathcal P+\mathcal X$, specify all
positive source Gram blocks of the complete parent.

\paragraph{An explicit list of 51 testing functions.}
Fix $P_*\ge1$, put $t=p/P_*$, and write
\[
 \phi=(X,Y,W,wW,wZ,pX)^T,\qquad
 y_1=t,\quad y_2=t(w/\delta)^2,\quad y_3=w/\delta.
\]
Take the $42$ labelled functions consisting of every entry of $\phi$
multiplied by one of
\[
 1,\quad\sqrt{1+y_1},\sqrt{1-y_1},\quad
 \sqrt{1+y_2},\sqrt{1-y_2},\quad
 \sqrt{1+y_3},\sqrt{1-y_3},
\]
and the nine further functions
\[
 pwZ,\ pY,\ pwY,\ wY,\ pwX,\ pW,\ pZ,\ Z,
                         \quad wZ\sqrt{1-t^2}.
\]
All radicands are nonnegative on $0\le p\le P_*$,
$0\le w\le\delta$. This list is independent of all source variables.
For the physical slack multiply its entries by $\kappa_t$.
For signed angles compose the list with the fixed angle-sign and
state-sign gauges. No state-coordinate inverse or sign-dependent
source factor occurs.

\paragraph{Polynomial weighted blocks.}
Here is the complete smaller weighted representation. A row
$(v,g,H)$ means $g\,v^THv$. Write $[v]=vv^T$, and put
\begin{gather*}
 t_B=\sqrt B,\quad t_H=\sqrt{H_0},\quad
 \ell_J=(t_B/q,-bt_B/q)^T,\quad \ell_G=(0,t_H/q)^T,\\
 h_s=\frac1{P_*}(q,-sz/q,-s(1-r)/q)^T\quad(s=1,4/3),\\
 \nu=1-a^2-c^2=\frac{Qz(1-r)}{H_0},\qquad \psi=(Y,wZ)^T.
\end{gather*}
Define source matrices
\begin{align*}
 G_{RS}&=(2-f/2)I_2,\\
 G_{EL}&=\tfrac32f[h_{4/3}]
               +2f\left(\frac{P_*^2}{1-r}-1\right)[h_1],\\
 G_{JG}&=\tfrac32f[\ell_J]+\tfrac56f[\ell_G]
                                      +2(1-f)[(1,-1)^T],\\
 G_{pJG}&=f\left(2-\frac2{3P_*^2}\right)[\ell_J]
       +\frac{2f}{3P_*^2}\{\nu([\ell_J]+[\ell_G])
                                      +[c\ell_J-a\ell_G]\}
       +2(1-f)[(1,0)^T].
\end{align*}
Their coefficients are nonnegative. In the order
$(R,S,J_*,\bar E,G_*)=\mathsf L\phi$, the source matrix is
\[
 \mathsf L=\begin{pmatrix}
 1&0&0&1&0&0\\0&0&0&0&0&1\\
 0&t_B/q&0&0&-bt_B/q&0\\
 0&0&q/P_*&0&0&0\\0&0&0&0&t_H/q&0
 \end{pmatrix}.
\]
Let $\Pi$ be the permutation with
$\Pi y=(R,S,J_*,\bar E,G_*)^T$, and define
$\widetilde T_i=\Pi T_i\Pi^{\mathsf T}$ from the three matrices
in the theorem. Let $M_i$ be their stated row-sum bounds.
Thus the seven main rows are
\[
 \left(\phi,1,\mathsf L^T(f/2-\rho)I_5\mathsf L\right),\qquad
 \left(\phi,1\pm y_i,
              \mathsf L^T\frac{M_iI_5\pm\widetilde T_i}{2}\mathsf L\right),
 \quad i=1,2,3.
\]
The other twelve rows are
\[
\begin{array}{c|c|c}
 v&g&H\\\hline
 (X+wW,pX)^T&1&G_{RS}\\
 (W,pY,pwZ)^T&1&G_{EL}\\
 \psi&1&G_{JG}\\
 wZ&1-t^2&2ft_H^2/(3q^2)\\
 \psi&p^2&G_{pJG}\\
 (pwX,pwY)^T&1&2I_2\\
 (W-pY,w(W-pY))^T&1&4\operatorname{diag}(C-fC_0,C)\\
 (Y,wY)^T&D_p&4\operatorname{diag}(A-fA_0,A)\\
 pW-pwX+Y-wZ&1&4m\\
 wW+X-pwY-pZ&1&4D\\
 pX+Z&D_w&4\beta\\
 Y-wZ&D_p&4(\alpha-f\alpha_0)
\end{array}
\]
Their sum is exactly \eqref{eq:allz-explicit-parent}. There are
$19$ blocks, of total order
\[
 7\cdot6+(2+3+2+1+2)+2+2+2+4=62,
\]
with maximum order six. All matrix entries are source-only. They
are explicit sums or congruences of nonnegative source rank-one
matrices, diagonal matrices, or the proved row-sum matrices.
The testing scalar multipliers are nonnegative polynomials.

To recover the stated scalar list, factor each source Gram in the
source field and split $D_p=1+p^2$, $D_w=1+w^2$, and $p^2$ into
polynomial testing squares. Every resulting ungated root is a
source-linear combination of the fourteen displayed monomials;
the only remaining radicals are the six main gate roots and
$\sqrt{1-t^2}$ multiplying $wZ$. Thus no additional testing function
is introduced by any of the retained full events.



\begin{corollary}[All Alice angles in the bounded chart]\label{cor:alice}
Choose $P_*=2^{24}$. On the native source domain
\[
 0<z\le1,\qquad 0<x,\sigma\le\frac{2^{-78}}{1+z},
\]
the union of this construction and the inner-angle certificate\footnote{\href{https://github.com/anubhav70747/q222-finite-sdp/tree/main/extracted-lifts/local-order669}{A fixed finite C1 testing family and an evaluated rational local pencil}, Theorem~1.} covers every projective Alice angle and
$|w|\le2^{-36}$, for every state and arbitrary $H$.
\end{corollary}
\begin{proof}
The present theorem covers $|u|\ge2^{-24}$, including infinity.
For $|u|\le2^{-24}$, the positive-$\lambda$ inner testing
condition $\sqrt{|u|},|w|\le2^{-12}$ holds. The nonpositive branch
has the weaker requirement $|u|,|w|\le2^{-12}$. The present source
box is contained in both inner source boxes. The cases meet at
$|u|=2^{-24}$ and include $u=0$ literally.
\end{proof}


\paragraph{The rational affine pencil.}
The weighted vectors in the table determine a finite semidefinite
representation without retaining the nonlinear source formulas as constraints.
Let $c=(d,m,C,A,D,\alpha,\beta)$, and let $G_b$ be unrestricted positive
semidefinite matrices of the displayed orders. Define the certificate cone by
\[
 \widehat F_c=\sum_{b=1}^{19} g_b v_b^{\mathsf T}G_bv_b,
 \qquad G_b\succeq0.
\]
The identity is an equality of polynomials in $p,w,W,X,Y,Z$.
Every $v_b$ and $g_b$ is fixed once $P_*$ is chosen. Their rational
coefficients make coefficient comparison a rational linear system
\[
 \mathcal A\,\gamma=\mathcal B\,c,
\]
where $\gamma$ lists the $176$ upper-triangular Gram entries.
For $P_*=2^{24}$ there are $80$ coefficient equations of rank $79$.
Eliminating the pivot entries leaves $97$ free entries $y$ and gives
the rational pencil $L_{62}(c,y)$ of total order $62$.
The accompanying \texttt{SOURCE\_PENCIL.coo.json} lists its $338$
nonzero affine coefficients.

The two directions of this representation follow from the same elimination.
A feasible pencil reconstructs positive semidefinite $G_b$ satisfying every
polynomial coefficient equation. Conversely, for any certificate, take its
$97$ free Gram entries as $y$; the pivot equations reconstruct the remaining
entries of those same matrices. Thus the pencil represents this weighted
certificate cone exactly.

Sharing $c$ with the inner-angle pencil gives the intersection
\[
 L_{731}(c,y,z)=\operatorname{diag}\bigl(L_{669}(c,y),L_{62}(c,z)\bigr).
\]
It has $139$ blocks and $2791$ auxiliary coordinates. Its source admission
and testing coverage are precisely those of Corollary~\ref{cor:alice}.
The file \texttt{ALL\_ALICE\_ORDER731.coo.json} gives all its entries.

\paragraph{Remaining domains.}
For $0<z\le R$ one sufficient admission box is
$x,\sigma\le2^{-30}P_*^{-2}/(1+R)$; the testing family is unchanged
as $R$ varies. Every germ with $x,\sigma$ infinitesimal and $z$
bounded satisfies \eqref{eq:domain}. Thus this is one testing family
for the inverse-Alice part of all those bounded-$z$ germs, without
choosing a real upper bound for $z$ first. It also covers unbounded
$z$ whenever $(1+z)x,(1+z)\sigma\le\eta$.
It does not cover a native source with larger weighted openings,
nor reciprocal Bob angles $|w|>\delta$. For $z>1$, the existing
inner extension has an $R$-dependent testing interval; the present
outer identity alone does not remove that separate inner issue.
Literal zero openings excluded by the native formulas, such as $z=0$, are handled
in their other source charts and are not obtained by division here.
No full-$Q_{222}$ adequacy statement follows.
\end{document}
